\documentclass[10pt]{article}
\usepackage[margin=20mm]{geometry}
\usepackage{amsmath,amssymb,mathtools,fontspec,unicode-math,xurl,hyperref}
\IfFontExistsTF{DejaVu Serif}{\setmainfont{DejaVu Serif}}{\setmainfont{Latin Modern Roman}}
\setmathfont{latinmodern-math.otf}
\setlength{\parindent}{0pt}\setlength{\parskip}{6pt}\sloppy
\title{Initial review of character extensions in type D}
\author{}
\date{26 October 2026}
\begin{document}
\maketitle
Abstract: This claim-focused review checks the equivariant lifting mechanism, elementary extension arguments, the signed-trace counting step, and the final rank-induction structure. No substantive defect is established in the inspected application to the main theorems. A distinction between fixed-point count and full orbit size in Proposition 6.4 remains an explicit scope question, not a proved counterexample. The substantial Lie-theoretic labelling and descent prerequisites have not all been independently audited. Scientific acceptance remains pending.\par
\section*{Identity, version, and coverage}
The bound work is Britta Späth, Extensions of characters in type D and the inductive McKay condition, II, Inventiones mathematicae 242 (2025), 45–122, DOI `10.1007/s00222-025-01354-9` [1]. The acquired publisher version of record has 78 PDF pages and 2,208,737 bytes. PDF page \(p\) corresponds to printed page \(44+p\); page references below use PDF numbering unless explicitly marked printed.\par
The actual text and image inspection covers pp. 1–9, 12–17, 23–24, 50–54, 68–69, and 73–77: 29 pages. The selected claims are Theorems A–C, Condition 2.1 and Hypothesis 2.2, Lemmas 3.2, 3.5–3.6, the architecture of Theorem 3.3 and Proposition 3.11, Theorem 3.19, the reduction in Proposition 6.4, the trace calculation in Proposition 6.6, Lemma 7.1 and the final use of Proposition 7.7. Reading a theorem’s dependency statement is not counted as verifying the dependency’s entire proof.\par
\section*{Definitions and exact claims}
Here \(G=\mathbf G^F\) is the specified universal cover of a simple group of Lie type, with regular embedding into \(\widetilde G\) and graph-field automorphism group \(E(G)\). The introductory restriction on exceptional universal covers is explicit (p. 3); their separate character-table treatment is not checked here. The diagonal action factors through the finite abelian group \[\mathcal Z_F=\widetilde G/(GZ(\widetilde G)).\] Condition \(A'(\infty)\) asserts an \(E(G)\)-stable diagonal transversal. Condition \(A(\infty)\) additionally asks each chosen representative \(\chi\) to extend to \(G E(G)_\chi\). It does not say that every diagonal conjugate extends to the same graph-field subgroup. That distinction matters in the two-character orbits of Proposition 7.7.\par
Theorem B asserts an equivariant bijection with conjugacy classes of pairs \((s,\phi)\), including unipotent characters of possibly disconnected finite centralizers. Orbit sums, irreducible characters, rational series and geometric series must remain distinct. Theorem C concerns every finite group and the prime \(3\), and uses both the new global condition and previously established local conditions. It is not an assertion here of all-prime McKay solely from Theorem A.\par
\section*{Checked equivariant and extension arguments}
\subsection*{Lifting orbit bijections}
Lemma 3.2 (p. 12) requires three separate inputs: stable transversals, equality of stabilizer subgroups in \(Z\), and an equivariant bijection of \(Z\)-orbit sets. Choose stable transversals \(T_I,T_J\) and the induced map \(j:T_I\to T_J\). Defining \[J(zx_0)=zj(x_0),\qquad z\in Z,\ x_0\in T_I,\] is well-defined because two choices of \(z\) differ by \(Z_{x_0}=Z_{j(x_0)}\). The same condition gives injectivity and surjectivity on each orbit. For \(e\in E\), \[J(e(zx_0))=J(e(z)\,ex_0)=e(z)j(ex_0)=e(z)e(j(x_0))=e(J(zx_0)).\] This verifies full \(ZE\)-equivariance. Equality of orbit sizes alone would not justify this construction.\par
In Theorem 3.19 (pp. 23–24), the stabilizers on both sides are identified with the same inverse image under \(\widehat\omega_s\) of an annihilator subgroup. This is the appropriate subgroup equality, not merely a cardinality comparison. Its input from Proposition 3.17(c) remains an unaudited dependency here. Proposition 3.11 (pp. 16–17) supplies the second stable transversal using an invariant extension from Corollary 3.9 and Clifford correspondence.\par
The Clifford step needs only a twist by a linear character trivial on \(C=C_{\mathbf H}^{\circ}(s)^F\). If two extensions on the inertia group differ by a character of its quotient by \(C\), that character extends to the abelian component quotient and induction gives the required twist. This spells out the standard reading sufficient for the orbit argument; no stronger kernel condition is needed.\par
\subsection*{Central kernels and wreath products}
Lemma 3.5 (pp. 13–14) uses \(N\cap C_G(N)=Z(N)\leq\ker\chi\). Therefore \[NC_G(N)/C_G(N)\cong N/Z(N)\] carries \(\chi\), and an extension through the image of \(G_\chi\) in \(A_\chi\) inflates to the desired extension. The central-kernel hypothesis is essential and is present.\par
For Lemma 3.6 (p. 15), conjugating the factors first makes characters in the same \(Y\)-orbit equal. The stabilizer splits along equal-character blocks. On one block, an extension represented on \(V\) gives the tensor product representation on \(V^{\otimes n}\); permutation of tensor factors implements \(\mathcal S_n\) and satisfies the semidirect-product relations. Its restriction to \(X^n\) is the specified irreducible tensor character, so it is an extension. The general case is the product over blocks. This checks the elementary wreath-product step without assuming extension from invariance alone.\par
Theorem 3.3 uses algebraic components and automorphisms arising from them, not an unrestricted automorphism assertion about arbitrary small finite groups. Malle’s extension theorem and the unipotent identifications invoked there are still external prerequisites, including small-field cases.\par
\section*{Counting and the two-extension alternative}
\subsection*{The signed trace}
Set \(N=|\operatorname{Irr}(G^F)^D|\), let \(b\) be the number of those characters without extension to \(G^F D\), and set \(M=|\operatorname{Irr}(G^{F_0})^{\langle\gamma\rangle}|\). In Proposition 6.6 (pp. 53–54), the chosen basis contributes \(+1\) for an extendible invariant character and \(-1\) for a nonextendible invariant character. A pair exchanged by \(\gamma\) has trace zero. Consequently \[\operatorname{Tr}(\gamma_1)=N-2b=M,\qquad
 b=\frac{N-M}{2},\qquad N-b=\frac{N+M}{2}.\] Both signs and the factor \(1/2\) are consistent. The equality of traces through Shintani descent is assumed from the cited construction; it was not computationally reproduced. Given a subset of nonextendible characters with cardinality \((N-M)/2\), finite-set inclusion and equal size prove that it is the complete exceptional set. Cardinality alone, without nonextendibility, would not suffice.\par
\subsection*{Lemma 7.1}
The quotient in Lemma 7.1 (pp. 68–69) is the stated abelian \(2\)-group with independent cyclic factors generated by the images of \(d,f,c\); the images of \(d,c\) are involutions. Let \(\varphi\) extend \(\chi\) to \(\langle X,f\rangle\). Since this quotient over \(X\) is cyclic, extensions differ by its linear characters, and conjugation by \(c,d\) acts trivially on those linear characters. The noninvariance hypothesis for extensions to \(\langle X,d\rangle\) forces \[\varphi^d=\varphi\nu,\qquad \nu^2=1,\quad \nu\neq1.\] Conjugation by \(c\) gives either \(\varphi^c=\varphi\) or \(\varphi^c=\varphi\nu\). Thus exactly one of \(c\) and \(dc\) fixes \(\varphi\). The corresponding next quotient is cyclic, so an invariant \(\varphi\) extends. Twisting \(\varphi\) cannot change either invariance test; hence the other subgroup admits no extension of \(\chi\). This last observation justifies the nonexistence assertion for all extensions, not only for a selected one.\par
An independent exact calculation tested both alternatives in groups of order 16. On tuples \((z,d,c,f)\in\mathbf F_2^4\), use \[x y=(z+Z+f(D+aC),\,d+D,\,c+C,\,f+F),\qquad a\in\{0,1\},\] where \(y=(Z,D,C,F)\). The central subgroup \(X=\langle z\rangle\) has character \(\chi(z)=-1\), and the quotient is \(C_2^3\). In both cases \([d,f]=z\), so \(f\) interchanges the two extensions to \(\langle X,d\rangle\). For \(a=0\), there are four linear extensions to \(\langle X,f,c\rangle\) and none to \(\langle X,f,dc\rangle\); for \(a=1\) the counts reverse. Associativity was checked on all 4096 triples in each group. These finite examples illustrate the lemma; the preceding analytic argument supplies the general check.\par
\section*{Scope question and main-theorem dependence}
Proposition 6.4(ii.2), p. 52 (printed p. 96), assumes \(|\operatorname{Irr}(\bar\chi)^D|=2\). Its last proof case explicitly adds \(|\operatorname{Irr}(\bar\chi)|=2\) and then uses the order-two diagonal stabilizer on p. 53. These are different predicates. This review has not established the remaining general case where the full orbit is larger but exactly two points are \(D\)-fixed, nor produced a character of the specified Lie-type groups refuting the proposition. The general wording is therefore an unresolved scope question, not a confirmed defect.\par
There is decisive counterevidence to inferring failure of Theorem A from that question: the application on p. 76 invokes Proposition 7.7 to obtain a full two-element orbit whenever a nonextendible character occurs, and otherwise every invariant character extends. Thus this application uses precisely the stronger alternative treated in the proof. No unsupported case is required at that interface, conditional on the cited construction of the exceptional sets.\par
The rank induction in the proof of Theorem A (p. 76) starts at \(l=4\), where the smaller-rank hypothesis is empty. The induction hypothesis covers every smaller rank and every field exponent required by Hypothesis 2.2; Theorem 6.1 supplies the untwisted case, followed by the twisted case through Theorem 5.1. Mentioning dependence between twisted and untwisted arguments does not alone establish a circular proof. The complete internal verification of those descent arguments remains outside this review.\par
For Theorem B, Theorem A supplies \(A'(\infty)\) and Theorem 3.19 supplies the equivariant bijection. For Theorem C (pp. 76–77), the text separately invokes the reduction theorem, the non-Lie-type results, and \(A(1),B(1),A(2),B(2)\). Those external results were not independently audited.\par
\section*{Disposition}
No substantive issue has been established in the inspected main-theorem application. No repair or author adoption is claimed. A limited title/correction search and the publisher record did not reveal an article-specific correction; this is not proof of absence. The associated arXiv record was identified, but no later author text was used to silently replace the bound journal version.\par
Item | Status\par
Issue judgment | No Issue Found within the stated scope; one unresolved proposition-scope question.\par
Verified reasoning | Orbit lift, elementary extension mechanisms, signed-trace algebra, two-alternative lemma, final conditional induction.\par
Repair/adoption | No substantive correction proposed; author adoption not requested.\par
Remaining work | Full component-group and descent proofs, exceptional-set construction, Proposition 6.4 in its general wording, external prerequisites.\par
Acceptance | Pending; human reviews empty.\par
\section*{Supporting checks}
Its actual output and execution receipt are in `evidence/independent\_checks.json` and `evidence/independent\_execution\_log.json`. These are reviewer-written finite-group checks; no author code or character-table software was executed. The source manifest contains the original PDF and publisher HTML, with no computational code or essential missing supplement identified.\par
Uninspected PDF pages are 10–11, 18–22, 25–49, 55–67, 70–72 and 78. In particular, most of Sections 4–6, the full proof of Proposition 7.2 and the complete proof of Theorem 3.21 are not certified by this report. On inspected pages, substantial cited Lie-theoretic inputs remain dependencies.\par
> 00004TOKEN\par
The reviewing assistant performed all scientific work, calculations, compilation and page inspection. The main dispatcher provides scheduling and attributed UI observations only. The environment self-designation is GPT-6 Astra Pro; backend-to-UI equivalence is not independently authenticated. No human review or acceptance has been entered.\par
1 B. Späth, Extensions of characters in type D and the inductive McKay condition, II, Invent. Math. 242 (2025), 45–122. DOI `10.1007/s00222-025-01354-9`. Bound publisher PDF; publisher record consulted 26 October 2026. References numbered in the body as paper dependencies refer to this article’s numbering, not independently consulted full texts.\par
\end{document}
